\chapter{The Options Market-Making Houses}\label{ch:in:the-options-market-making-houses}

In the early 1980s a group of friends from college began trading on the floor of the Philadelphia Stock Exchange, and
in 1987 they joined together in one firm. Two years later two traders working side by side on the floor of the
Amsterdam equity options exchange founded another. In 2002 a third firm took its place in the index-options pit on the
floor of the Chicago Board Options Exchange. Each firm tells its own origin this way, on its own pages, and each is
today an options market maker that quotes on screens around the world. The floor taught them to price every strike,
to hold the book, and to hedge what they held; they still hire and train people to do the same, now at the speed of
software. This chapter describes these houses as employers: where they come from, how they train, how their work
differs from that of the electronic firms of chapter~2, and what the one of them that publishes its results shows.

\section{Three lineages: Amsterdam, Chicago, Philadelphia}

Listed options were once traded on exchange floors, among them in Chicago, Amsterdam and Philadelphia. The market
makers of those floors were small partnerships of \omterm{def:in:proprietary-market-makers-and-high-frequency-firms:floor}{floor traders} (chapter~2) who priced options with
models and hedged them with the underlying shares. Some of their firms became the options houses of today.

\begin{definition}[Options market-making house]\label{def:in:the-options-market-making-houses:house}
An \emph{options market-making house}\index{options market-making house} is a \omterm{def:in:a-map-of-the-industry:ptf}{principal trading firm} whose core business
is quoting exchange-listed options, across many strikes and expiries of many underlyings, usually under market-maker
appointments with quoting obligations, and hedging the resulting positions in the underlying and in other options.
\end{definition}

The definition describes a firm; the role it plays on a venue, the options market maker with its quoting obligations
and its protections, is Book~1's (chapter~24). \Cref{fig:in:the-options-market-making-houses:lineage} draws the origins
that the houses state. Three of them tie their founding to a specific floor; the others give a city and a year.

\begin{omfigure}
\begin{tikzpicture}[font=\footnotesize,>=stealth,
  venue/.style={draw,rounded corners,fill=omDef!12,minimum width=3.6cm,minimum height=0.6cm,align=center},
  firm/.style={draw,rounded corners,fill=omProp!8,minimum width=3.2cm,minimum height=0.5cm,align=center}]
\node[venue] (ams) at (0,0) {Amsterdam equity\\options exchange};
\node[venue] (phl) at (4.4,0) {Philadelphia\\Stock Exchange};
\node[venue] (cbo) at (8.8,0) {Chicago Board\\Options Exchange};
\node[firm] (opt) at (-1.0,-1.5) {Optiver (1986)};
\node[firm] (imc) at (-1.0,-2.3) {IMC (1989)};
\node[firm] (sig) at (4.4,-1.5) {Susquehanna (1987)};
\node[firm] (gro) at (4.4,-3.1) {Group One (1989)};
\node[font=\scriptsize,gray] at (4.4,-3.65) {San Francisco};
\node[firm] (pea) at (9.8,-1.5) {PEAK6 (1997)};
\node[firm] (bel) at (9.8,-2.3) {Belvedere (2002)};
\node[font=\scriptsize,gray] at (-1.0,-2.85) {Amsterdam};
\node[font=\scriptsize,gray] at (9.8,-2.85) {Chicago};
\draw[->,thick,omDef] (imc.east) -- ++(0.7,0) |- (ams.south east);
\draw[->,thick,omDef] (sig.north) -- (phl.south);
\draw[->,thick,omDef] (bel.west) -- ++(-0.6,0) |- (cbo.south west);
\end{tikzpicture}

\omcaption{Where six options market makers say they began: a solid arrow means that the firm's own page ties its
founding to trading on that floor; firms without an arrow state only a city and a year. Sources: the firms' pages
(ledger F1, F3--F5, F7, F8).}\label{fig:in:the-options-market-making-houses:lineage}
\end{omfigure}

Four of the six firms give founding years in the 1980s, the decade of the floors; one each in the 1990s and
2000s. Two of the six began in Amsterdam; chapter~27 describes the tax and visa arrangements that bring staff to
the cities where the houses are.

\begin{definition}[Spin-out]\label{def:in:the-options-market-making-houses:spinout}
A \emph{spin-out}\index{spin-out} is a firm founded by people who leave an existing firm to do the same or a related
business on their own. Unlike a spin-off, the parent firm does not distribute or keep the new firm's shares.
\end{definition}

Lineages in this industry are often told as chains of \omterm{def:in:the-options-market-making-houses:spinout}{spin-outs}, from one house to the traders who left it. None of
the six firms states such an origin on its own pages, so the chapter's graph has no \omterm{def:in:the-options-market-making-houses:spinout}{spin-out} edges: a lineage built
from public sources is thinner than the one the industry tells about itself, and chapter~28 returns to what the public
record does show about people leaving to start firms.

\section{Training programmes as the firms describe them}

An options house turns graduates into traders who price and hedge thousands of instruments at once, and it says so
when it recruits. One house, whose founders built their first strategies with ``poker experience'', states that ``we use strategy games and decision science to
teach many of our new employees''; its campus page carries an intern's account of a ten-week internship spent on an
index volatility desk. Another describes its education as ``learning built for live markets and real-time systems''.
The content taught is the content of Books~1, 5 and 11: option payoffs and parity, the Greeks, volatility surfaces,
inventory and hedging. What the houses add is practice: simulated markets, trading games, and months of supervised
work before a trader holds risk alone.

\begin{remark}[What a programme is]\label{rem:in:the-options-market-making-houses:programme}
A training programme is the firm's investment in turning a hire into someone it trusts with risk. Its length, its
content and whether a hire who leaves early must repay anything are terms of the employment contract, and differ by
firm and by country. The firms' pages describe the experience, not the terms; a candidate reads the contract
(chapter~13).
\end{remark}

\section{How an options house differs from an electronic delta-one firm}

Chapter~2's firms mostly quote instruments with one price each: a share, a future, a currency pair. An options house
quotes a surface. The differences follow.

\begin{enumerate}
\item \textbf{The number of instruments.} One underlying with twenty listed expiries and fifty strikes each has
$20\times50\times2=2\,000$ options series. A house quoting a few hundred underlyings maintains hundreds of thousands of
quotes, updated together as the underlying moves (Book~11, chapter~19).
\item \textbf{The risk held.} A delta-one market maker aims to be flat within minutes. An options house holds a book of
volatility, skew and time-decay exposures for days or weeks, hedging the direction (delta) continuously and the rest
by trading other options.
\item \textbf{The obligations.} Market-maker appointments on options exchanges carry quoting obligations over many series
(Book~1, chapter~24): the house must be present when it would rather not be.
\item \textbf{The model.} Prices come from a volatility surface fitted live to the market (Book~11, chapter~19; Book~5),
not from a single fair price; a trader's judgement acts on the surface's parameters.
\end{enumerate}

For an employee the consequence is a different job. A trader at an options house manages a book of Greeks and decides
where the surface should be; a trader at a delta-one firm supervises strategies that race for a single price. Both
use the same computers; the questions are different.

\begin{example}[One house's scale]\label{ex:in:the-options-market-making-houses:scale}
One house, founded on the Amsterdam floor, states that it trades on 120 venues with 2\,070 employees in ten offices.
If each venue lists options on only a hundred underlyings of 2\,000 series, the house could quote on the order of
$120\times100\times2\,000=24$ million series; the actual number is smaller, because many venues list the same
underlyings and many series are never quoted, but the order of magnitude explains why quoting is software and a
trader's work is supervision and pricing judgement.
\end{example}

\section{Size and footprint}

Unlike most private trading firms, one options house publishes its results every year, with a short annual review.

\begin{dated}{2026-09}{The options houses in their own words}\label{dat:in:the-options-market-making-houses:houses}
\textbf{Optiver} (results for 2025, published 31 March 2026): net trading income \euro4\,556 million (2024: \euro3\,494
million; 2023: \euro2\,773 million); net profit attributable to equity holders \euro1\,769 million (2024: \euro1\,369
million; 2023: \euro1\,158 million); total equity \euro5\,490 million (2024: \euro4\,905 million); ``more than 2\,000
employees'' (2\,100 in the 2024 release), established in Amsterdam in 1986, with offices in ten other cities.
\textbf{IMC}: founded in 1989; 2\,070 employees across ten offices, trading on 120 venues.
\textbf{Group One}: founded in San Francisco in 1989; ``over 100'' people in four cities; an options market maker or
specialist on every major US exchange group. \textbf{Akuna Capital}: ``a leading proprietary trading firm specializing
in options market making''. The first two describe their options business in their own words: one prices ``1M+''
instruments ``from options and ETFs to equities, bonds and currencies''; the other calls itself ``a top-3 liquidity
provider by volume traded in listed options globally''.
\end{dated}

The published results make one house's economics visible (\cref{fig:in:the-options-market-making-houses:results}).
Net trading income rose by 26\% in 2024 and by 30\% in 2025; the net profit margin on it was 39.2\% and 38.8\%; the
return on average equity was 30.4\% and 34.0\%. With 2\,100 employees in 2024, net trading income was \euro1.66 million
per head; with ``more than 2\,000'' in 2025, at most \euro2.28 million per head. Chapter~12 puts these numbers beside
the other firms that file.

\begin{omfigure}
\begin{tikzpicture}
\begin{axis}[width=10.5cm,height=5.2cm,ybar,area legend,bar width=14pt,xmin=2022.4,xmax=2025.6,ymin=0,ymax=6,
  xtick={2023,2024,2025},xticklabel style={/pgf/number format/1000 sep={}},ylabel={\small \euro\ billion},
  tick label style={font=\footnotesize},grid=major,grid style={gray!20},table/col sep=comma,
  legend style={legend cell align=left,font=\scriptsize,at={(0.5,-0.25)},anchor=north,/tikz/every even column/.append style={column sep=8pt}},legend columns=3]
\addplot[fill=omDef!55,draw=omDef] table[x=year,y=nti]{figdata/industry/03-the-options-market-making-houses/results.csv};
\addplot[fill=omThm!55,draw=omThm] table[x=year,y=profit]{figdata/industry/03-the-options-market-making-houses/results.csv};
\addplot[fill=gray!40,draw=gray] table[x=year,y=equity]{figdata/industry/03-the-options-market-making-houses/results.csv};
\legend{net trading income,net profit,total equity (year end)}
\end{axis}
\end{tikzpicture}

\omcaption{One privately owned options house's published results, 2023--2025: net trading income rose from \euro2.8
billion to \euro4.6 billion and net profit from \euro1.2 billion to \euro1.8 billion. Data: Optiver's results releases
for 2024 and 2025, through \texttt{in\_options.results}.}\label{fig:in:the-options-market-making-houses:results}
\end{omfigure}

\begin{remark}[Why one house publishes]\label{rem:in:the-options-market-making-houses:why}
A privately owned firm that publishes results chooses to: for its counterparties, its lenders and its recruits. The
numbers are the firm's own, audited but unfiled in any register a reader can check line by line, and the review is
written by the firm. They are a primary source for what the firm reports, not for how it compares with firms that
report nothing.
\end{remark}

\section{What the filings show for the options houses}

Three of the sourced market makers, Optiver, IMC and Susquehanna, are the options houses of this chapter, and together
they file enough labour condition applications (chapter~14) to be compared with the other market makers without
publishing any one firm's pay.

\begin{dated}{2025-09}{Options houses and other market makers in the filings}\label{dat:in:the-options-market-making-houses:lca}
Fiscal 2025 labour condition applications, offered base: the three options houses, 130 applications, median \$150\,000
(10th--90th percentile \$84\,800--200\,000); the other seven market makers, 237 applications, median \$200\,000
(\$130\,000--300\,000). Researchers \$150\,000 against \$200\,000; software engineers \$146\,100 against \$190\,000;
traders \$200\,000 against \$192\,500. At wage level~I: 23 of the options houses' 130 applications, 8 of the others' 237.
Base salary only.
\end{dated}

\begin{omfigure}
\begin{tikzpicture}
\begin{axis}[width=10.5cm,height=6cm,xmin=50,xmax=320,ymin=0.4,ymax=8.6,ytick=data,
  yticklabels from table={figdata/industry/03-the-options-market-making-houses/options_pay.csv}{label},
  yticklabel style={font=\scriptsize},xticklabel style={font=\footnotesize},xlabel={\small offered base, \$ thousand},
  grid=major,grid style={gray!20},table/col sep=comma]
\addplot[draw=none,mark=none,error bars/.cd,x dir=both,x explicit,error mark=none,error bar style={omDef,line width=1pt}]
  table[x=p50,y=pos,x error plus expr=\thisrow{p90}-\thisrow{p50},x error minus expr=\thisrow{p50}-\thisrow{p10}]
  {figdata/industry/03-the-options-market-making-houses/options_pay.csv};
\addplot[draw=none,mark=none,error bars/.cd,x dir=both,x explicit,error mark=none,error bar style={omDef!40,line width=6pt}]
  table[x=p50,y=pos,x error plus expr=\thisrow{p75}-\thisrow{p50},x error minus expr=\thisrow{p50}-\thisrow{p25}]
  {figdata/industry/03-the-options-market-making-houses/options_pay.csv};
\addplot[only marks,mark=|,mark size=4.5pt,line width=1.5pt,omThm] table[x=p50,y=pos]{figdata/industry/03-the-options-market-making-houses/options_pay.csv};
\end{axis}
\end{tikzpicture}

\omcaption{Offered base in fiscal 2025 labour condition applications: the three options houses against the other
market makers, by role family: 10th to 90th percentile (thin), interquartile range (thick), median (mark). Pooled over
three or more employers; no firm's figures are shown. Data: \texttt{data/industry/lca\_options.csv}, through
\texttt{in\_options.options\_pay}.}\label{fig:in:the-options-market-making-houses:lca}
\end{omfigure}

The options houses' median offer is \$50\,000 lower (95\% bootstrap interval $-\$53\,900$ to $-\$35\,000$) for all titles
and for researchers, and their software engineers' \$43\,900 lower; their traders' is level with the others'
(\cref{fig:in:the-options-market-making-houses:lca}). Two cautions come before any reading. The options houses file
many more applications at the entry level, 17.7\% at wage level~I against 3.4\%, consistent with the graduate
training programmes described above; and base salary is only part of pay at a market maker, where chapter~14 found the
implied ratio of total pay to base largest. The filings show where base offers sit, not what the jobs pay.

\section{Tutorial: the houses' origins and one house's numbers}\label{tut:in:the-options-market-making-houses:tut}

\begin{tutorial}
\textbf{Goal.} Build the lineage graph from the firms' own statements, and compute one house's economics from its
published results. \textbf{End state:} \cref{fig:in:the-options-market-making-houses:lineage}'s edges and the ratios of
the section above.
\begin{tutsteps}
\item \textbf{The graph.} \texttt{data/industry/lineage.csv} lists three venues and six firms; three rows carry an edge
from a firm to the floor its founders traded on. \texttt{in\_options.graph()} loads it with \texttt{firm.lineage},
which refuses a node or an edge without a ledger row
(\cref{lst:in:the-options-market-making-houses:add}).
\omcode{code/firm/lineage/firm_lineage.py}{51}{61}{An edge enters the graph only with a source and between known
nodes.}\label{lst:in:the-options-market-making-houses:add}
\item \textbf{Query.} \texttt{g.by\_place()} groups the firms by city, \texttt{g.by\_decade()} counts founding decades
(1980s 4, 1990s 1, 2000s 1), and \texttt{g.unsourced()} lists the three firms whose origin the public record
used does not tie to a floor.
\item \textbf{The numbers.} \texttt{in\_options.metrics()} computes growth, margin, return on average equity and income
per head from \texttt{data/industry/optiver\_results.csv}, treating ``more than 2\,000 employees'' as a lower bound on
headcount, and so an upper bound on income per head.
\end{tutsteps}
\end{tutorial}

\textbf{What to change next.} Add the founding year and origin of another options house from its own pages, and see
whether it changes the decade counts; compute what 2025 income per head would be if headcount were 2\,500
(exercise~7).

\section{Build: the lineage graph}\label{bld:in:the-options-market-making-houses:lineage}

\begin{build}
\bfield{Purpose} A sourced genealogy of trading firms, for this chapter's floors and for chapter~28's question of where
new firms come from.
\bfield{Interface} \texttt{firm.lineage}: \texttt{Node(name, kind, year, place, ledger)}, \texttt{Edge(src, dst,
relation, ledger)}; \texttt{Graph.load(path)}, \texttt{add\_node}, \texttt{add\_edge}; \texttt{firms}, \texttt{by\_place},
\texttt{by\_decade}, \texttt{children}, \texttt{ancestors}, \texttt{unsourced}.
\bfield{Rules} No node or edge without a ledger row; an edge joins two known nodes; a founding year is optional and is
never inferred.
\bfield{Acceptance tests} \texttt{code/firm/lineage/tests/}: grouping by place and decade, children and ancestors on a
constructed graph; unsourced and dangling edges are refused.
\bfield{Stretch} Export the graph to TikZ with computed positions; add acquisition edges from filings (a listed firm's
10-K names its acquisitions).
\end{build}

\begin{omsources}
\item The firms' own pages: IMC (About us), Susquehanna (About, Campus), Belvedere Trading, PEAK6, Group One, Akuna
Capital, Optiver (home page, Education).
\item Optiver, results releases for 2024 (27 March 2025) and 2025 (31 March 2026).
\item Cboe Global Markets, Form 10-K for 2025, on the Cboe Options trading floor.
\end{omsources}

\section{Exercises}

\begin{exercise}[$\star$]\label{exo:in:the-options-market-making-houses:1}
How many options series does an underlying with 15 expiries and 40 strikes each list, counting calls and puts?
\end{exercise}

\begin{exercise}[$\star$]\label{exo:in:the-options-market-making-houses:2}
From the dated box, compute the house's net profit margin on net trading income in 2023.
\end{exercise}

\begin{exercise}[$\star$]\label{exo:in:the-options-market-making-houses:3}
Which of the six firms of \cref{fig:in:the-options-market-making-houses:lineage} tie their founding to a floor, and
which floor?
\end{exercise}

\begin{exercise}[$\star\star$]\label{exo:in:the-options-market-making-houses:4}
Compute the house's growth in net profit in 2024 and 2025, and explain why profit grew more slowly than net trading
income in 2024.
\end{exercise}

\begin{exercise}[$\star\star$]\label{exo:in:the-options-market-making-houses:5}
The house's net trading income per head was \euro1.66 million in 2024. What range does 2025 income per head lie in if
headcount was between 2\,000 and 2\,400?
\end{exercise}

\begin{exercise}[$\star\star$]\label{exo:in:the-options-market-making-houses:6}
Give two reasons an options house holds risk longer than an electronic delta-one market maker.
\end{exercise}

\begin{exercise}[$\star\star\star$]\label{exo:in:the-options-market-making-houses:7}
\emph{Coding.} Compute the return on average equity for 2024 and 2025 with \texttt{in\_options.metrics}, and the 2025
income per head for a headcount of 2\,500.
\end{exercise}

\begin{exercise}[$\star\star\star$]\label{exo:in:the-options-market-making-houses:8}
\emph{Find the flaw.} ``Options houses earn about \euro2 million per employee: the published results of one house prove
it for the industry.''
\end{exercise}

\section{Problem: Three Floors, One Screen}

\begin{problem}[{Weekend problem --- three floors, one screen}]\label{pb:in:the-options-market-making-houses:1}
A graduate choosing between an options house and a delta-one trading firm wants to understand the houses' origins,
their work and one house's economics.

\medskip\noindent\textbf{Part I --- Origins.}
\begin{enumerate}
\item Define an \omterm{def:in:the-options-market-making-houses:house}{options market-making house}, and say how it differs from Book~1's options market maker.
\item Name the three floors of \cref{fig:in:the-options-market-making-houses:lineage} and the firm tied to each.
\item Count the six firms' founding decades.
\item Define a \omterm{def:in:the-options-market-making-houses:spinout}{spin-out} and say why the graph has none.
\item What does the Chicago options exchange's floor still do, according to its operator?
\end{enumerate}

\medskip\noindent\textbf{Part II --- The work.}
\begin{enumerate}[resume]
\item How many series does an underlying with 20 expiries and 50 strikes list?
\item Give the four differences between an options house and a delta-one firm.
\item What does a training programme teach, by the firms' account, and what does it add to the books?
\item Why is quoting software and a trader's job supervision and pricing judgement?
\item What does a candidate read to learn a programme's terms?
\end{enumerate}

\medskip\noindent\textbf{Part III --- One house's numbers.}
\begin{enumerate}[resume]
\item Give net trading income and net profit for 2023--2025.
\item Compute growth in net trading income in 2024 and 2025.
\item Compute the net profit margin in each year.
\item Compute the return on average equity in 2024 and 2025.
\item Compute net trading income per head in 2024, and its upper bound in 2025.
\end{enumerate}

\medskip\noindent\textbf{Part IV --- The verdict.}
\begin{enumerate}[resume]
\item State the \emph{named result}: the house's net trading income, margin and return on average equity in 2025, and
the share of the six firms tied to a floor by their own account.
\item Why is ``more than 2\,000'' an upper bound on income per head?
\item Why can one house's published numbers not stand for the industry's?
\item Which public source would tell the graduate most about a house that does not publish?
\item In two sentences, what should the graduate take from the chapter to the choice?
\end{enumerate}
\end{problem}

\section{Interview questions}

\begin{interviewq}[$\star$ \iqroles{trader}]\label{iq:in:the-options-market-making-houses:1}
Why does an options market maker hold risk overnight when an equity market maker tries to end the day flat?
\end{interviewq}

\begin{interviewq}[$\star$ \iqroles{trader, researcher}]\label{iq:in:the-options-market-making-houses:2}
A single stock has options at 12 expiries and 30 strikes. How many quotes does a market maker maintain if it quotes
both sides of every series?
\end{interviewq}

\begin{interviewq}[$\star\star$ \iqroles{trader}]\label{iq:in:the-options-market-making-houses:3}
An exchange's market-maker appointment obliges you to quote a wide range of series. Why would a firm accept the
obligation, and what does it cost?
\end{interviewq}

\begin{interviewq}[$\star\star$ \iqroles{researcher, risk}]\label{iq:in:the-options-market-making-houses:4}
A firm's net profit rose 29\% and its equity rose 12\%. What happened to its return on equity, roughly, and what would
you check before concluding the business got better?
\end{interviewq}

\begin{interviewq}[$\star\star$ \iqroles{trader}]\label{iq:in:the-options-market-making-houses:5}
What would you expect to learn in the first six months at an options house that a textbook does not teach?
\end{interviewq}

\begin{interviewq}[$\star\star\star$ \iqroles{trader, risk}]\label{iq:in:the-options-market-making-houses:6}
Your options book is long gamma and short vega after a busy day. Describe what you would do before the close, and why.
\end{interviewq}
